
\section*{Reproducibility}

Every numerical claim that enters the proof of Theorems A--D is traceable to a
machine-readable certificate in the accompanying verification package, save for
the few premises named below that are established by hand or by inspection
instead, and one that is not checked at all: the correctness of the interval
arithmetic itself, which no computer-assisted proof in this field checks. The
article also prints numbers that carry no proof, and each is marked as such
where it stands --- at the number, or at the paragraph or section that holds it
--- and \textbf{that labelling, and not any list, is what the
reader should go by}. They fall into two kinds, and the instances named below
are the principal ones rather than an exhaustive roll.
\emph{Not rigorous:} quantities
computed in floating point and labelled so on the spot --- the eigenvalue
$\lambda_{2}$ as a function of $\ell$ (\S\ref{sec:7.4}), the exploratory scan
of the constants over the segment (\S\ref{sec:7}), the diagnostics of the
routes that failed (\S\ref{sec:4.2}), $\mathrm{cond}(M')$ at $\ell = 5/2$ and
the diagnostic finite-difference agreement quoted at $\ell = 5/2$ (both
\S\ref{sec:5.3.3}), and the
cross-assembly of $\DT$ by finite differences at three interior points
(\S\ref{sec:3.9}).
\ifx\JOURNALVERSION\undefined
The side observation on the family $\alpha^{-k\ell}$ and the one-off
experiments on how far the floor in the model self-check can be pushed down are
of the same kind and are labelled in the same way.
\else
Two further quantities of this kind --- the side observation on the family
$\alpha^{-k\ell}$ (\S\ref{sec:8}) and the one-off experiments on how far the
floor in the model self-check can be pushed down (\S\ref{sec:3.10}) --- are not
printed in this article at all: the labelling rule applies to them where they
stand, in that document.
\fi
The labelling rule is about quantities computed from the object under study.
Three kinds of printed number fall outside it and are not marked individually:
reported run times; inventory counts --- counts of files, of stored records or
of sampled points, and dataset sizes; and the parameters and coordinates of the
computation itself --- working precisions, truncation orders, numbers of nodes
and arcs, file and line references, and dates. A wall clock reading, like a
count of files or a working precision, is a description of how the work was
done, never a claim about the fixed point, and none of them appears in any
inequality. The counts that do carry statements --- the 46 existence boxes, the
706 tube boxes --- are not of this kind and are certified like everything else.
\emph{Rigorous but not recomputed from the package:} products that are
cross-checked against the certificate values rather than rebuilt by the
acceptance run --- \notrecomputedroll{}

One qualification belongs here rather than only at its own place. The two
elementary conditions (E1)--(E2) of \S\ref{sec:10.2} --- $|c'| < r'$ and
$|1-c'| < r'$ --- are premises of the lemmas behind Theorem D, and although the
parameters $(c', r')$ they constrain come from the certificates, the
inequalities themselves are not gated by the \texttt{one\_box} certifier; they
are checked separately, for every deposited layer and cell, by the numbers gate
(\texttt{tests/\allowbreak \_\allowbreak numbers\_\allowbreak check.py}, block K4:
$40$ objects, worst case $(0.70, 0.75)$ with margin $0.05$).
Every other premise of Theorems A--D that is meant to be machine-checked is
gated automatically, with no exception. The box \texttt{t0695}, whose
intersection with $[2,4]$ degenerates to a point, was one until the containment
check stopped clipping to $[2,4]$ before comparing: $[4.0000, 4.0012]
\subset \texttt{gtop} = [3.9775, 4.0075]$, so the term \texttt{extra} is
justified on it as well (premise a7). Nothing in the
statement about $[2,4]$ rests on it, and \S\ref{sec:4.8} says so at the place
where the covering is closed. The premises that are proved analytically rather
than gated are marked \texttt{[P]} in
Table~\ref{tab:10.11-proof}, the map of the proof, and the one
premise that is neither --- the correctness of the interval arithmetic itself
(premise A4) --- is listed in Table~\ref{tab:10.11-assumptions}, the registry
of assumptions, as common to the whole field and not checked
here.
One premise outside Theorems A--D is in the same position: the
single-valuedness of the branch $U_G^{5/2}$ over the Kantorovich ball at
$\ell = 5/2$ rests on an order-of-magnitude estimate rather than on a gate, as
\appref{sec:5} states at the point of use. It is one premise but it stands in
two places --- the existence certificate (\texttt{kant\_\allowbreak l25.json})
and the hyperbolicity certificate (\texttt{eigen\_\allowbreak l25.json},
\texttt{hyperb\_\allowbreak l25.json}), neither of which gates
$\Real U$ --- and both are independent pointwise verifications that carry no
step of Theorems A--D.
The package carries the
engine, the certificates it produced, the tests that falsify the majorants, and
a SHA-256 manifest of all 960 files. Two inputs are deliberately not in it, and
each is identified at the point where it is used: the march bank behind the
independent cross-check of \S\ref{sec:3.8}, and the driver script of
the acceptance run, \texttt{verify\_\allowbreak pkg.sh}, which is supplied to the
referee separately --- the individual gate scripts that it calls are in the
package. All
dependencies are open source; the interval computations use the Arb ball
arithmetic \cite{johansson2017arb} through \texttt{python-flint} 0.9.0
\cite{pythonflint} (FLINT 3.6.0 \cite{flint}).

A reviewer who wants to check the work without reading the code in full can do
so in five steps, in this order; the fifth is two separate runs, so six
programs of the package --- besides the external check named in step 1, which
the package cannot perform on itself:

\begin{enumerate}
\item \texttt{python \_\allowbreak manifest.py verify} --- the package matches its manifest.
  The verifier is itself part of the manifest and therefore cannot certify
  itself; check at least once with an external tool, for instance
  \texttt{sha256sum -c -{}-quiet MANIFEST.sha256}.
\item \texttt{python 04-\allowbreak CODE/\allowbreak tests/\allowbreak \_\allowbreak classic\_\allowbreak constants.py} --- the engine reproduces
  the published Feigenbaum values $\delta$ and $\alpha$ at the endpoints of the
  interval (15 digits for $\delta$ and 14 for $\alpha$ at $\ell = 2$). This is a check external to the
  derivation of the derivative.
\item \texttt{python 04-\allowbreak CODE/\allowbreak tests/\allowbreak \_\allowbreak tau\_\allowbreak test.py} --- the tail majorants really are
  majorants: the test is built so that the pre-correction version must fail.
\item \texttt{python 04-\allowbreak CODE/\allowbreak cert.py} on one existence box --- for instance
  the anchor \texttt{ganchor} ($\ell_{0} = 3$) --- checked against
  \texttt{03-\allowbreak CERTIFICATES/\allowbreak existence/\allowbreak }. This script is the
  command-line runner over the existence engine that the article cites by file
  and line as \texttt{rig2fix3.py}: it imports that module and is not a second
  implementation.
\item \texttt{python 04-\allowbreak CODE/\allowbreak tube\_\allowbreak run.py} on one lane of the plan, checked
  against \texttt{03-\allowbreak CERTIFICATES/\allowbreak tube/\allowbreak }, and
  \texttt{04-\allowbreak CODE/\allowbreak tube\_\allowbreak accept.py} for the full acceptance run.
  \texttt{tube\_\allowbreak run.py} is the lane driver: it certifies the boxes of one lane
  through \texttt{hyp\_\allowbreak rouche.py} and the tube engine \texttt{tube\_\allowbreak lib.py}, both
  of which the article cites by name at the places where they are used.
  \texttt{tube\_\allowbreak accept.py} builds no
  certificate of its own: it reads the box files, recomputes the full test from
  their fields in interval arithmetic, and applies the acceptance conditions of
  \S\ref{sec:4.8} to the whole chain.
\end{enumerate}

Expected numbers for each step are listed in the reproduction guide inside the
package.

\textbf{Which of them are reproducible, and in what sense.} Steps 1--3 read the
deposited artefacts and recompute verdicts from their fields; they are deterministic
for a fixed library stack and must give exactly the printed numbers. Steps 4 and 5 are
not of that kind: they \emph{regenerate} a certificate, and the float seed
\texttt{solve\_\allowbreak fixed\_\allowbreak point} behind $Y0$, $\tau_{F}$,
$\rho^{\star}_{b}$ and $\mathrm{extra}$ is not deterministic across library builds
(\S\ref{sec:3.6}, \S\ref{sec:3.7}: $Y0$ varies by tens of percent between runs, and on an
independent regeneration the degenerate box \texttt{gtop} was observed at
$\mathrm{extra} = 2.96\cdot 10^{-6}$, above the maximum printed here). A regeneration
that returns different values of these four quantities is therefore not a failure to
reproduce. What must hold on a regeneration is the verdict, not the digits: the
Newton--Kantorovich gate passes at the same $\rho'$, $N_p$ and $K$, and the resulting
ball still lies inside the uniqueness ball $1.75\cdot 10^{-7}$. The constants printed in
Theorems A--C are upper bounds attained by the \emph{deposited} run; they are held to the
manifest, not to any future one.

\paragraph{What the acceptance run does and does not establish.} The driver
\texttt{verify\_\allowbreak pkg.sh} performs \textbf{26} checks --- the package
tests, the acceptance runs and the external cross-checks --- and a clean run
reports all 26 passing. That statement should be read for exactly what it is:
the verdict scripts of the package pass on the packaged artefacts. The run
reads the 752 box certificates and the layer certificates as data and
recomputes the tests from their fields; it does not regenerate them. Steps 4
and 5 of the route above are where a reader regenerates one existence box and
one lane of the tube.
Three things it cannot establish by its internal-consistency gates alone ---
the principal ones, not a closed list. The limits already named above in this
section are limits of the same run and are not repeated here: the premises
(E1)--(E2), checked by the numbers gate rather than by the certifier; the
premises proved analytically and marked \texttt{[P]}; and the products that are
cross-checked against the certificates rather than rebuilt by the run.
One further limit belongs here, because the article leans on the mechanism
five times: the three checks that read prose ---
\texttt{\_\allowbreak numbers\_\allowbreak check.py},
\texttt{\_\allowbreak refs\_\allowbreak check.py} and
\texttt{\_\allowbreak appendix\_\allowbreak check.py} --- read the package's
own copy of the article, \texttt{01-\allowbreak ARTICLE.md}, and not this
manuscript --- a Russian-language copy --- and no gate of the package opens a
\TeX{} file at all. What they
certify is that the numbers, addresses and field names of that copy agree with
the artefacts. The manuscript is held to the same numbers by hand, and its
agreement with the package copy is not machine-checked. Two
checks of the run reach past those gates, in different ways: one compares
against data external to the package
(\texttt{\_\allowbreak classic\_\allowbreak constants.py}), the other steps
outside the certificates' own fields by rebuilding $\DT$ along a second,
independent route
(\texttt{\_\allowbreak dt\_\allowbreak fd\_\allowbreak cross.py}). It is these
two, not the count of passing gates, that bear on the second item below:

\begin{itemize}
\item \textbf{The correctness of the interval arithmetic itself.} The gates
  that carry the proof verify their inequalities \emph{under} the ball
  arithmetic of \texttt{python-flint}/\texttt{arb}, and a systematic error in
  its rounding would turn them green on a false result. Not every step of the
  run is of that kind, and those that are not say so where they stand ---
  among them the
  finite-difference cross-assembly compares
  a difference whose truncation is not rigorously enclosed against a fixed
  threshold (\S\ref{sec:3.9}); the \texttt{kmode} step measures a distance in
  float64 after casting down from \texttt{arb} (\S\ref{sec:9}). None of the
  latter carries a claim of Theorems A--D. The risk here is premise A4 of the
  registry in \S\ref{sec:10.11}, external to this work and common to the field.
\item \textbf{A coordinated error shared by the model and the engine.} If the
  code that assembles the centre $P(\ell)$ or the operator value carried a
  systematic shift, both the certificate and the check that recomputes it from
  the certificate's own fields would move together, and the gates would pass;
  \S\ref{sec:4.10} makes the same point about a gate not seeing what makes its
  own quantity smaller. The two checks named above bear on it, and neither
  closes it. The agreement with the classical constants of the literature is
  decisive but lives \textbf{at the two ends of the segment} only ---
  \S\ref{sec:3.9} says of it, in those words, ``only at the ends, along the
  leading mode, and no more than that''. The
  finite-difference assembly of $\DT$ is run \emph{by the driver} at
  \textbf{one} interior point, $\ell = 3$, the two others ($\ell = 5/2$ and
  $\ell = 3.9$) being reproduced by hand with a command-line argument
  (\S\ref{sec:3.9}); so no automatically run check assembles $\DT$ a second
  time at a \textbf{non-integer} interior $\ell$, by a route independent of
  the analytic derivative. And --- as \S\ref{sec:3.9} states in
  those words --- it is independent of the derivation of the \emph{derivative},
  but not of the engine as a whole: it calls the same operator-value code
  \texttt{compute\_\allowbreak F\_\allowbreak poly}. An error in that code
  would therefore be shared by the certificates, by the tube, by the anchor and
  by this cross-check alike. For interior $\ell$ that is the weakest link in
  the whole construction, and we state it here rather than let the count of
  passing gates suggest that something closes it.
\item \textbf{The provenance of the artefacts.} \texttt{\_manifest.py} checks
  the bytes against its own list and is itself listed in that manifest, so it
  cannot certify itself; step~1 of the route above therefore calls an external
  \texttt{sha256sum} as well. Nor can the run establish that the certificates
  and the scripts are an independent realisation of the mathematics stated
  here: they are one realisation, checked against itself.
\end{itemize}

\noindent A passing run is thus a necessary condition and a convenience, not a
second proof. The reader who wants an independent handle should start from the
two ends of the segment, where the values are known from the literature
independently --- to $325$ digits at $\ell = 4$ (BOT, Table~II) and to
thousands at $\ell = 2$, with which our own values agree to fourteen and
fifteen digits respectively --- and from \S\ref{sec:3.9}.

\section*{Data availability}

All source code, interval certificates, verification scripts and the SHA-256
manifest supporting the results of this article are openly available in a Zenodo
archive --- with the exception named above, the driver script
\texttt{verify\_\allowbreak pkg.sh} of the acceptance run, every gate script it calls being
in the archive --- version DOI: \texttt{10.5281/\allowbreak zenodo.\allowbreak 22306911} (all versions: \texttt{10.5281/\allowbreak zenodo.\allowbreak 22306910}).
That script, its instructions
(\texttt{README-verify.md}) and the retroactive check of the logarithm tail
(\texttt{logtail\_\allowbreak retro.py}) are \emph{not part of the package}
either: all three are
submitted to the editorial office together with
this manuscript, as separate files, and are to be passed to the referee with
it. Their SHA-256 digests are
\begin{flushleft}\footnotesize\ttfamily
verify\_pkg.sh\\
a52ec3ca8f891835ca4a27b500916f44f9610f7480f14fa12f79e32caf090766\\[2pt]
README-verify.md\\
a024110ca18047a4c44a313081a55fd3b50e8ea58e757f8c82a7f4cf8113bafc\\[2pt]
logtail\_retro.py\\
ce695a4e3a042cc2e1395d747dddf46a746cea32b1f03523ac71e9fd4546d855
\end{flushleft}
\noindent The version DOI is the one to follow: the proof depends on the exact file hashes, and the concept DOI resolves to whatever version is newest. Every numerical claim that enters the proof of
Theorems A--D is traceable to a machine-readable certificate indexed in that
archive; the Reproducibility section lists the numbers that are not.
Exploratory and one-off computations were carried out in the course of this
work whose outputs are not deposited: the recomputation of $|\lambda_{2}|$ and
$|\lambda_{3}|$ beyond $\ell = 4$ (\S\ref{sec:7.4}), the bank of march
certificates behind the independent cross-check of \S\ref{sec:3.8}, and a
number of smaller diagnostic runs. \textbf{Each is identified as undeposited at
the point where its result is quoted}; that labelling, and not any list, is
what the reader should go by. No step of any proof in this article depends on
data outside the archive.

\ifx\JOURNALVERSION\undefined
The material that the journal version carries as separate supplementary
sections --- what the certificates do not establish, the exploratory scan of
the constants over $\ell$, the numerical side observation on the family
$\alpha^{-k\ell}$, and the layer diagnostics --- is presented inline in this
preprint, and is deposited in the same Zenodo archive.
\else
The supplementary material referred to in the text (sections numbered
$\mathrm{S}n$) is a separate document: it states what the certificates do not
establish, reports the exploratory float scan of the constants over $\ell$,
describes the numerical side observation on the family $\alpha^{-k\ell}$, and
collects the layer diagnostics. It
carries no step of any proof in this article --- every step of every proof,
including the derivation of the $\tau$ bounds behind Theorem A and the
certification protocols behind Theorem D, is in the article itself, in the
appendices. The supplementary material is submitted alongside the manuscript,
is included in full in the preprint version, and is deposited in the same
Zenodo archive.
\fi

\section*{Declarations}

\subsection*{Funding}
This research did not receive any specific grant from funding agencies in the
public, commercial, or not-for-profit sectors.

\subsection*{Competing interests}
The author declares no competing interests.

\subsection*{Author contributions}
Sole author: conceptualization, proofs, software, writing.

\subsection*{Acknowledgements}
The author thanks Fredrik Johansson and the \texttt{python-flint} developers for
the Arb ball arithmetic library on which every rigorous bound in this work
rests.
